\section{Variable Selection and Coefficient Diagnostics}

This section presents the full OLS regression diagnostics for all four
candidate CDX IG fair-value models estimated on the training set (first 80\% of
the sample). For each model we report: coefficient estimates, standard errors,
\(t\)-statistics, and model-level information criteria (AIC, BIC) that penalise
over-parameterisation.

\subsection*{Model-level summary}

Table~\ref{tab:model_summary} compares the four candidate specifications on
key fit and selection metrics. Lower AIC and BIC values indicate a more
parsimonious fit.

\begin{table}[H]
\centering
\small
\caption{Model-level diagnostics (training set).}
\label{tab:model_summary}
\begin{tabular}{lrrrrrr}
\toprule
Model & \(n\) & \(k\) & \(R^2\) & Adj.\ \(R^2\) & AIC & BIC \\
\midrule
CrossMarket\_d       & 916 & 6 & 0.516 & 0.514 & 4\,015 & 4\,044 \\
RatesVolEquity\_d    & 952 & 6 & 0.485 & 0.483 & 4\,201 & 4\,231 \\
FXBasisPolicy\_d     & 667 & 6 & 0.057 & 0.049 & 3\,548 & 3\,575 \\
InflationMacro\_d    & 940 & 5 & 0.146 & 0.142 & 4\,632 & 4\,657 \\
\bottomrule
\end{tabular}
\end{table}

\textbf{Observations:}
\begin{itemize}
  \item On in-sample \(R^2\) alone, CrossMarket\_d ranks highest (0.516).
    However, this model includes iTraxx IG as a regressor---an index that is
    economically near-collinear with CDX IG---making it partially tautological.
    It is therefore not selected as the primary fair-value specification despite
    its in-sample fit.
  \item RatesVolEquity\_d achieves the highest \emph{out-of-sample} \(R^2\)
    (0.698, as reported in the main document) despite a lower in-sample \(R^2\),
    indicating better generalisation and less overfitting.
  \item AIC and BIC are not directly comparable across models with different
    sample sizes (FXBasisPolicy\_d uses a shorter series due to ESTR data
    availability). Among the three models sharing comparable sample sizes
    (\(\sim\!920\)--\(950\)), CrossMarket\_d has the lowest AIC/BIC, followed
    by RatesVolEquity\_d, with InflationMacro\_d substantially higher.
  \item FXBasisPolicy\_d has very low explanatory power (\(R^2 = 0.057\)),
    confirming that FX and short-rate variables alone are insufficient to
    explain IG credit spread movements.
\end{itemize}

\subsection*{Coefficient detail: RatesVolEquity\_d (selected model)}

Table~\ref{tab:coef_rve} presents the coefficient estimates for the selected
specification.

\begin{table}[H]
\centering
\small
\caption{OLS coefficients --- RatesVolEquity\_d (training set, \(n = 952\)).}
\label{tab:coef_rve}
\begin{tabular}{lrrrr}
\toprule
Variable & Coefficient & Std.\ Error & \(t\)-statistic & Sig.\ (5\%) \\
\midrule
Intercept      &  $+0.077$ & $0.071$ &  $+1.08$ & No  \\
\(\Delta\)VIX  &  $+0.268$ & $0.053$ &  $+5.06$ & Yes \\
\(\Delta\)S\&P &  $-0.036$ & $0.003$ & $-11.87$ & Yes \\
\(\Delta\)UST10 & $-4.127$ & $2.074$ &  $-1.99$ & Yes \\
\(\Delta\)USYC &  $+1.059$ & $2.970$ &  $+0.36$ & No  \\
\(\Delta\)SOFR &  $+0.234$ & $0.503$ &  $+0.46$ & No  \\
\bottomrule
\end{tabular}
\end{table}

\textbf{Interpretation:}
\begin{itemize}
  \item \(\Delta VIX\) and \(\Delta S\&P\) are highly significant
    (\(|t| > 5\)). Higher volatility widens spreads; equity appreciation
    tightens them. These are the dominant drivers.
  \item \(\Delta UST10\) is marginally significant at the 5\% level
    (\(t = -1.99\)). The negative sign indicates that rising yields are
    associated with tighter spreads---consistent with a growth-driven regime
    where higher rates reflect improving fundamentals.
  \item \(\Delta USYC\) and \(\Delta SOFR\) are not individually significant,
    suggesting that once the 10-year yield level and VIX are controlled for,
    the additional rates information from curve slope and overnight rate
    contributes limited marginal explanatory power. These variables may
    nonetheless improve out-of-sample stability by capturing secondary rate
    dynamics.
\end{itemize}

\subsection*{Coefficient detail: CrossMarket\_d}

\begin{table}[H]
\centering
\small
\caption{OLS coefficients --- CrossMarket\_d (training set, \(n = 916\)).}
\label{tab:coef_cm}
\begin{tabular}{lrrrr}
\toprule
Variable & Coefficient & Std.\ Error & \(t\)-statistic & Sig.\ (5\%) \\
\midrule
Intercept           & $-0.022$ & $0.071$ &  $-0.31$ & No  \\
\(\Delta\)ITRX IG   & $+0.824$ & $0.063$ & $+13.03$ & Yes \\
\(\Delta\)ITRX XO   & $+0.028$ & $0.015$ &  $+1.88$ & No  \\
\(\Delta\)ESTX600   & $+0.081$ & $0.025$ &  $+3.23$ & Yes \\
\(\Delta\)EURUSD    & $-17.86$ & $15.41$ &  $-1.16$ & No  \\
\(\Delta\)DEM10     & $+1.431$ & $2.139$ &  $+0.67$ & No  \\
\bottomrule
\end{tabular}
\end{table}

\textbf{Interpretation:}
iTraxx IG dominates (\(t = +13.03\)), which is expected given that it is the
European equivalent of CDX IG. The high coefficient (0.82) and significance
confirm near-collinearity between US and European IG CDS indices, making this
model unsuitable as an independent fair-value tool. ESTX600 is the only other
significant variable, consistent with the broader equity--credit channel.

\subsection*{Coefficient detail: InflationMacro\_d}

\begin{table}[H]
\centering
\small
\caption{OLS coefficients --- InflationMacro\_d (training set, \(n = 940\)).}
\label{tab:coef_im}
\begin{tabular}{lrrrr}
\toprule
Variable & Coefficient & Std.\ Error & \(t\)-statistic & Sig.\ (5\%) \\
\midrule
Intercept         & $+0.035$ & $0.093$ & $+0.38$ & No  \\
\(\Delta\)CO      & $-0.037$ & $0.114$ & $-0.33$ & No  \\
\(\Delta\)US5BEI  & $-22.62$ & $2.592$ & $-8.73$ & Yes \\
\(\Delta\)UST2    & $-13.11$ & $2.900$ & $-4.52$ & Yes \\
\(\Delta\)DEM2    & $-3.368$ & $4.443$ & $-0.76$ & No  \\
\bottomrule
\end{tabular}
\end{table}

\textbf{Interpretation:}
Following the exclusion of FR5BEI (see Section~2), this model now uses a
comparable sample size (\(n = 940\)) to the other specifications.
US breakeven inflation (\(t = -8.73\)) and the 2-year Treasury yield
(\(t = -4.52\)) are the only significant variables. The negative US5BEI
coefficient suggests that rising inflation expectations tighten IG
spreads---consistent with an environment where higher inflation proxies for
stronger nominal growth. The overall \(R^2\) of 0.146 remains low, indicating
that inflation and short-rate variables alone capture only a modest share of
CDX IG spread variation.

\subsection*{Coefficient detail: FXBasisPolicy\_d}

\begin{table}[H]
\centering
\small
\caption{OLS coefficients --- FXBasisPolicy\_d (training set, \(n = 667\)).}
\label{tab:coef_fbp}
\begin{tabular}{lrrrr}
\toprule
Variable & Coefficient & Std.\ Error & \(t\)-statistic & Sig.\ (5\%) \\
\midrule
Intercept             & $+0.007$ & $0.133$ & $+0.05$ & No  \\
\(\Delta\)EURUSD      & $-68.45$ & $26.57$ & $-2.58$ & Yes \\
\(\Delta\)EURUSD3mBS  & $+0.150$ & $0.029$ & $+5.15$ & Yes \\
\(\Delta\)SOFR        & $-6.527$ & $2.143$ & $-3.04$ & Yes \\
\(\Delta\)ESTR        & $-6.289$ & $6.860$ & $-0.92$ & No  \\
\(\Delta\)DEYC        & $-3.590$ & $4.372$ & $-0.82$ & No  \\
\bottomrule
\end{tabular}
\end{table}

\textbf{Interpretation:}
Three variables are individually significant (EURUSD, EURUSD3mBS, SOFR), yet
the overall model \(R^2\) is only 0.057. This indicates that while FX and
funding-cost variables have statistically detectable relationships with CDX IG,
they explain very little of its total variance. The FX basis swap coefficient
is positive: a wider basis (more expensive USD funding) is associated with
wider IG spreads, consistent with global dollar-funding stress raising credit
risk premia.

\subsection*{Summary of variable selection}

\begin{table}[H]
\centering
\small
\caption{Significance summary across all candidate models (5\% level).}
\label{tab:sig_summary}
\begin{tabular}{lccccc}
\toprule
 & RatesVol & CrossMkt & InflMacro & FXPolicy & Significant in \\
Variable & Equity & & & & \# models \\
\midrule
VIX        & \checkmark &            &            &            & 1 \\
S\&P       & \checkmark &            &            &            & 1 \\
UST10      & \checkmark &            &            &            & 1 \\
ITRX IG    &            & \checkmark &            &            & 1 \\
ESTX600    &            & \checkmark &            &            & 1 \\
US5BEI     &            &            & \checkmark &            & 1 \\
UST2       &            &            & \checkmark &            & 1 \\
EURUSD     &            &            &            & \checkmark & 1 \\
EURUSD3mBS &            &            &            & \checkmark & 1 \\
SOFR       &            &            &            & \checkmark & 1 \\
\bottomrule
\end{tabular}
\end{table}

The significance pattern confirms that each model captures a distinct channel:
risk appetite/volatility (RatesVolEquity), cross-region credit
(CrossMarket), inflation/short rates (InflationMacro), and
FX/funding (FXBasisPolicy). The selected model (RatesVolEquity\_d)
combines the two most economically interpretable and significant
drivers---equity risk appetite and implied volatility---with rates variables
that provide structural context, yielding the best out-of-sample performance.
